
\documentclass{article}
\usepackage{colortbl}
\usepackage{makecell}
\usepackage{multirow}
\usepackage{supertabular}

\begin{document}

\newcounter{utterance}

\centering \large Interaction Transcript for game `eqbench', experiment `eq\_bench', episode 110 with reflexion{-}Qwen3.5{-}9B.
\vspace{24pt}

{ \footnotesize  \setcounter{utterance}{1}
\setlength{\tabcolsep}{0pt}
\begin{supertabular}{c@{$\;$}|p{.15\linewidth}@{}p{.15\linewidth}p{.15\linewidth}p{.15\linewidth}p{.15\linewidth}p{.15\linewidth}}
    \# & \multicolumn{2}{c}{Player} && \multicolumn{2}{c}{Game Master} \\
    \hline

    \theutterance \stepcounter{utterance}  
    & & & \multicolumn{4}{p{0.6\linewidth}}{
        \cellcolor[rgb]{0.9,0.9,0.9}{
            \makecell[{{p{\linewidth}}}]{
                \texttt{\tiny{[P1$\langle$GM]}}
                \texttt{You have attempted this same task before and did not succeed. These are the plans you wrote after those attempts.} \\
\texttt{Trial \#1: New plan: I failed because I assigned a high score to "Provoked" while ignoring the nuance of Charlie's specific reaction to being called "bitter" and having their life described as "falling apart." While Charlie is certainly defensive, the dialogue escalates to personal attacks on Charlie's character and life situation, which would logically induce feelings of hurt, defensiveness, or perhaps even shame/embarrassment regarding the personal nature of the insults, rather than just general provocation. I should have recognized that "Embarrassed" is a strong candidate because Charlie is being personally attacked about their life falling apart, not just their professional critique. Next time, I will re{-}evaluate the emotional spectrum: acknowledge that while "Provoked" is high (9), "Embarrassed" should also be significant (e.g., 7 or 8) due to the personal nature of the insults, and ensure the scores reflect the intensity of the personal attack rather than just the argumentative tone. I will also verify that the sum of scores reflects the complexity of the situation, ensuring at least one score is >0 as required.} \\
\\ 
\texttt{Trial \#2: New plan: I failed because I assigned a high score to "Provoked" while ignoring the nuance of Charlie's specific reaction to being called "bitter" and having their life described as "falling apart." While Charlie is certainly defensive, the dialogue escalates to personal attacks on Charlie's character and life situation, which would logically induce feelings of hurt, defensiveness, or perhaps even shame/embarrassment regarding the personal nature of the insults, rather than just general provocation. I should have recognized that "Embarrassed" is a strong candidate because Charlie is being personally attacked about their life falling apart, not just their professional critique. Next time, I will re{-}evaluate the emotional spectrum: acknowledge that while "Provoked" is high (9), "Embarrassed" should also be significant (e.g., 7 or 8) due to the personal nature of the insults, and ensure the scores reflect the intensity of the personal attack rather than just the argumentative tone. I will also verify that the sum of scores reflects the complexity of the situation, ensuring at least one score is >0 as required.} \\
\\ 
\texttt{Trial \#3: New plan: I failed because I assigned a high score to "Provoked" while underestimating the impact of the personal attacks on Charlie's life situation ("falling apart") and character ("bitter"). While Charlie is certainly defensive, the dialogue escalates to insults about their personal life, which logically induces feelings of hurt, shame, or embarrassment regarding the vulnerability of their situation, rather than just general provocation. I should have recognized that "Embarrassed" is a strong candidate because Charlie is being personally attacked about their life falling apart, not just their professional critique. Next time, I will re{-}evaluate the emotional spectrum: acknowledge that while "Provoked" remains high (9), "Embarrassed" should also be significant (e.g., 7 or 8) due to the personal nature of the insults, and ensure the scores reflect the intensity of the personal attack rather than just the argumentative tone. I will also verify that the sum of scores reflects the complexity of the situation, ensuring at least one score is >0 as required.} \\
\\ 
\\ 
\\ 
\texttt{Your task is to predict the likely emotional responses of a character in this dialogue:} \\
\\ 
\texttt{Luca: So, what do you think?} \\
\texttt{Charlie: It's... interesting.} \\
\texttt{Luca: Interesting? That's it?} \\
\texttt{Charlie: Well, it's not exactly a compliment, is it?} \\
\texttt{Luca: You haven't even tasted it yet. } \\
\texttt{Charlie: I don't need to taste it to know it's a mess.} \\
\texttt{Luca: A mess? You're calling my masterpiece a mess?} \\
\texttt{Charlie: If the shoe fits...} \\
\texttt{Luca: You're just a critic, Charlie. You don't know what it takes to create something.} \\
\texttt{Charlie: And you don't know what it takes to appreciate something. } \\
\texttt{Luca: I've been a chef for over 20 years. I've earned my stars.} \\
\texttt{Charlie: And I've been a critic for 15. I've earned my stripes.} \\
\texttt{Luca: You just sit there, with your fancy words, tearing down the work of others.} \\
\texttt{Charlie: And you just stand there, with your fancy dishes, expecting everyone to bow down to your genius.} \\
\texttt{Luca: You're just bitter because your life is falling apart.} \\
\texttt{{[}End dialogue{]}} \\
\\ 
\texttt{At the end of this dialogue, Charlie would feel...} \\
\texttt{Embarrassed} \\
\texttt{Provoked} \\
\texttt{Indifferent} \\
\texttt{Grateful} \\
\\ 
\texttt{Give each of these possible emotions a score from 0{-}10 for the relative intensity that they are likely to be feeling each.} \\
\\ 
\texttt{You must output in the following format, including headings (of course, you should give your own scores), with no additional commentary:} \\
\\ 
\texttt{Embarrassed: <score>} \\
\texttt{Provoked: <score>} \\
\texttt{Indifferent: <score>} \\
\texttt{Grateful: <score>} \\
\\ 
\\ 
\texttt{{[}End of answer{]}} \\
\\ 
\texttt{Remember: zero is a valid score, meaning they are likely not feeling that emotion. You must score at least one emotion > 0.} \\
\\ 
\texttt{Your answer:} \\
            }
        }
    }
     \\ \\

    \theutterance \stepcounter{utterance}  
    & \multicolumn{4}{p{0.6\linewidth}}{
        \cellcolor[rgb]{0.561,0.737,0.561}{
            \makecell[{{p{\linewidth}}}]{
                \texttt{\tiny{[P1$\rangle$GM]}}
                \texttt{Embarrassed: 2} \\
\texttt{Provoked: 9} \\
\texttt{Indifferent: 0} \\
\texttt{Grateful: 0} \\
            }
        }
    }
    & & \\ \\

    \theutterance \stepcounter{utterance}  
    & & & \multicolumn{2}{p{0.3\linewidth}}{
        \cellcolor[rgb]{0.9,0.9,0.9}{
            \makecell[{{p{\linewidth}}}]{
                \texttt{\tiny{[GM$|$GM]}}
                \texttt{Embarrassed: 2} \\
\texttt{Provoked: 9} \\
\texttt{Indifferent: 0} \\
\texttt{Grateful: 0} \\
            }
        }
    }
    & & \\ \\

    \theutterance \stepcounter{utterance}  
    & & & \multicolumn{2}{p{0.3\linewidth}}{
        \cellcolor[rgb]{0.9,0.9,0.9}{
            \makecell[{{p{\linewidth}}}]{
                \texttt{\tiny{[GM$|$GM]}}
                \texttt{Embarrassed: 2} \\
\texttt{Provoked: 8} \\
\texttt{Indifferent: 0} \\
\texttt{Grateful: 0} \\
            }
        }
    }
    & & \\ \\

    \theutterance \stepcounter{utterance}  
    & & & \multicolumn{2}{p{0.3\linewidth}}{
        \cellcolor[rgb]{0.9,0.9,0.9}{
            \makecell[{{p{\linewidth}}}]{
                \texttt{\tiny{[GM$|$GM]}}
                \texttt{game\_result = LOSE} \\
            }
        }
    }
    & & \\ \\

    \theutterance \stepcounter{utterance}  
    & & \multicolumn{4}{p{0.6\linewidth}}{
        \cellcolor[rgb]{0.95,0.95,0.95}{
            \makecell[{{p{\linewidth}}}]{
                \texttt{\tiny{[P1$\leftrightarrow$P]}}
                \texttt{Post{-}episode Reflexion (not sent to GM):} \\
\texttt{New plan: I failed because I assigned a high score to "Provoked" while underestimating the specific emotional impact of Luca's personal attacks on Charlie's life ("falling apart") and character ("bitter"). While Charlie is certainly defensive, the dialogue escalates from professional critique to exposing personal vulnerability, which logically induces feelings of shame, hurt, or embarrassment rather than just general provocation. I should have recognized that "Embarrassed" is a strong candidate because Charlie is being personally attacked about their life situation, not just their professional critique. Next time, I will re{-}evaluate the emotional spectrum: acknowledge that while "Provoked" remains high (9), "Embarrassed" should also be significant (e.g., 7 or 8) due to the personal nature of the insults, and ensure the scores reflect the intensity of the personal attack rather than just the argumentative tone. I will also verify that the sum of scores reflects the complexity of the situation, ensuring at least one score is >0 as required.} \\
            }
        }
    }
    & \\ \\

\end{supertabular}
}

\end{document}
